\hypertarget{transfer210-carrier-transfer-and-obstructions-for-fermat-hodge-packets}{%
\chapter{Carrier transfer and obstructions for Fermat Hodge
packets}\label{transfer210-carrier-transfer-and-obstructions-for-fermat-hodge-packets}}

\textbf{Komorebi P. --- research block v0.2.0}\\
\textbf{Status:} manuscript-level mathematical arguments with exact
computational certificates. The geometric argument is not formally
verified or peer reviewed. Exact certificates were independently
replayed for this edition. The degree-210 algebraicity conclusion uses
the degree-70 carrier proposition from v0.1.0 and explicitly identified
classical results. No general Hodge-conjecture solution or historical
novelty claim is made.

\hypertarget{transfer210-main-results-and-proof-boundaries}{%
\section{Main results and proof
boundaries}\label{transfer210-main-results-and-proof-boundaries}}

Let

\[
X_m^n=\{x_0^m+\cdots+x_{n+1}^m=0\}\subset\mathbf P^{n+1}_{\mathbf C}.
\]

This block obtains three different kinds of results.

\begin{enumerate}
\def\labelenumi{\arabic{enumi}.}
\tightlist
\item
  \textbf{Algebraicity transfer.} The rational primitive Hodge packet
  associated to \[
  b=(2,9,129,142,168,180)\pmod {210}
  \] is algebraic, using the v0.1.0 degree-70 carrier proposition and
  classical Fermat cycle correspondences. Its rational dimension is 48.
  Its full conductor is 210, so it is not a coordinate-power inflation
  of a lower-conductor packet. A 14-term signed identity supplies the
  transfer.
\item
  \textbf{Global matrix obstruction.} No nonnegative integral matrix
  capturing this packet can have positive normalized weights whose sum
  exceeds one. This is proved for every matrix size-six candidate of the
  previous criterion, not just for a bounded search. An exact dual
  certificate has two independent finite verifications, with a uniform
  argument covering all larger columns.
\item
  \textbf{Structural obstructions.} An infinite balanced family also
  cannot have such a matrix. A separate holomorphic-character-power
  criterion excludes all direct diagonal-quotient uniruled carriers for
  six explicitly checked test packets.
\end{enumerate}

The positive transfer and negative matrix result are compatible. They
prove that the direct matrix method is incomplete even on packets that
can be made algebraic by combining a geometric seed with classical cycle
operations.

The new standalone verifier checks 46 finite assertions. The preserved
predecessor verifier was separately rerun and passed all 24 of its
checks. Neither count is a formal verification of the geometric
theorems.

\hypertarget{transfer210-character-conventions}{%
\section{Character conventions}\label{transfer210-character-conventions}}

A full-support character is an ordered tuple \(a=(a_0,\ldots,a_{n+1})\)
of nonzero residues modulo \(m\), with sum zero modulo \(m\). Its
eigenspace in primitive cohomology is a complex line. For even \(n\),
its rational Galois packet is of middle Hodge type precisely when

\begin{equation}\label{eq:transfer210-2.1}
\sum_i\langle ta_i\rangle_m=\frac{n+2}{2}m
\quad\text{for every }t\in(\mathbf Z/m\mathbf Z)^\times.\end{equation}

Here residues are least positive representatives. This is the classical
Fermat character decomposition, see \cite{DaSilvaFermat}, Theorem 2.1.
Rational orbit dimension means the number of distinct \textbf{ordered}
unit conjugates, not the number of representatives after coordinate
permutations.

Write \(\nu(a)\) for its integer residue-multiplicity vector. Write
\(a\uplus b\) for concatenation regarded as a multiset. Permuting
coordinates is an algebraic automorphism, so multiset identities suffice
for algebraicity transfer. Write \(-a\) for coordinatewise negation
modulo \(m\).

A character is called algebraic when its eigenspace is contained in the
complex span of algebraic cycle classes. The rational orbit then lies in
the rational algebraic-cycle subspace: this subspace and the orbit
projector are rational, and the projector is a rational linear
combination of graphs of diagonal automorphisms.

\hypertarget{transfer210-classical-inputs-stated-separately-from-computation}{%
\section{Classical inputs, stated separately from
computation}\label{transfer210-classical-inputs-stated-separately-from-computation}}

We use the following established operations.

\textbf{Concatenation.} Algebraic balanced characters concatenate to an
algebraic character on the higher-dimensional Fermat variety.

\textbf{Complementary-pair removal.} If a character becomes algebraic
after concatenating any finite collection of pairs \((r,m-r)\), then it
was algebraic already.

These are the Shioda--Ran results recorded as Aoki \cite{Aoki}, Theorem
1-4(i),(ii). The equivariant Fermat join construction is also described
in \cite{DaSilvaFermat}, Theorem 2.2 and Corollary 2.3. The removal
operation is a geometric correspondence theorem, not a formal assumption
that arbitrary tensor factors can be cancelled.

\textbf{Four-entry blocks.} Every rational balanced four-entry packet is
algebraic, by the Lefschetz (1,1) theorem for the corresponding Fermat
surface.

\textbf{Standard blocks.} For an odd prime \(p\mid m\), put \(d=m/p\).
Aoki \cite{Aoki}, Theorem 2-1 constructs an algebraic cycle representing

\[
\sigma_{p,i}=(i,i+d,\ldots,i+(p-1)d,m-pi)
\]

under the relevant nonzero-entry and coprimality conditions. We only use
\((m,p,i)=(35,5,5)\) and \((35,7,2)\). Here \(\gcd(5,7)=\gcd(2,5)=1\),
and all entries are nonzero.

\textbf{Power transport.} The finite coordinate-power morphism
\(X_{km}^n\to X_m^n\) transports algebraic cycle classes by pullback.
The character is multiplied by \(k\).

\hypertarget{transfer210-source-access-boundary}{%
\subsection{Source-access boundary}\label{transfer210-source-access-boundary}}

The October audit obtained the original Aoki article. Theorem 1-4 states
concatenation and removal of complementary-pair auxiliary blocks.
Theorem 2-1 supplies the standard blocks below under its actual
nonzero-entry and coprimality conditions. The September source-access
limitation is preserved in the evidence archive and resolved for these
specific leaves in this edition. Neither theorem allows cancellation of
an arbitrary motive. The Shioda--Ran original proofs were not
independently re-audited.

\hypertarget{transfer210-the-inherited-geometric-seed}{%
\section{The inherited geometric
seed}\label{transfer210-the-inherited-geometric-seed}}

The authoritative predecessor is the included, unchanged archive

\texttt{predecessor/Komorebi\_Hodge\_Carrier70\_v0.1.0.zip}.

Its proposition asserts algebraicity of the 24-dimensional rational
packet

\begin{equation}\label{eq:transfer210-4.1}
a=(1,20,24,42,61,62)\pmod {70}.\end{equation}

The proof uses a concrete monomial correspondence to a normal weighted
degree-30 hypersurface with weights \((10,6,3,3,3,8)\). Its resolution
is shown to be uniruled, so the Conte--Murre Hodge theorem for uniruled
fourfolds applies. See \cite{ArapuraUniruled}, Corollary 4.2. Boundary
vanishing, finite étale descent and lowest-weight lifting then transfer
cycles back.

The predecessor's singular-locus, matrix and character checks have been
replayed, see \texttt{results/predecessor\_verification.json}. This does
not replace review of the predecessor's mixed-Hodge and birational
arguments.

Pullback under \(X_{210}^4\to X_{70}^4\) gives the algebraic source
packet

\begin{equation}\label{eq:transfer210-4.2}
s=3a=(3,60,72,126,183,186)\pmod {210}.\end{equation}

The source has conductor 70 and 24 distinct ordered unit conjugates. The
new target

\begin{equation}\label{eq:transfer210-4.3}
b=(2,9,129,142,168,180)\pmod {210}\end{equation}

has conductor 210 and 48 distinct ordered unit conjugates. Every
conjugate of both tuples sums to 630. Their Galois orbits are disjoint.
These are exact finite checks.

The target also occurs in the even-degree residual table accompanying
\cite{Jumagulov2026}. We use that table to select a research target, not
as proof that the entire census is correct or that this exact packet was
previously unknown in all literature.

\begin{figure}[tbp]\centering\includegraphics[width=0.98\linewidth]{23_transfer210_orbits.pdf}\caption{Source and target ordered unit orbits have dimensions 24 and 48, with conductors 70 and 210.}\end{figure}

\hypertarget{transfer210-an-exact-transfer-identity}{%
\section{An exact transfer identity}\label{transfer210-an-exact-transfer-identity}}

All entries in this section are modulo 210. Define

\[
P_k=(k,210-k)
\]

and

\[
\begin{aligned}
Q_3&=(3,105,108,204),\\
Q_9&=(9,105,114,192),\\
Q_{24}&=(24,105,129,162),\\
Q_{27}&=(27,105,132,156),\\
R_2&=(2,72,142,204).
\end{aligned}
\]

Every one of these four-entry blocks is balanced, as the exact verifier
checks against every unit modulo 210. They are algebraic by the surface
case.

Define the two longer blocks

\[
S_5=(30,60,72,114,156,198),
\qquad
S_7=(12,42,72,102,126,132,162,192).
\]

Their deflations are

\[
\begin{aligned}
S_5/6&=(5,10,12,19,26,33)=\sigma_{5,5}\quad\text{at degree }35,\\
S_7/6&=(2,7,12,17,21,22,27,32)=\sigma_{7,2}\quad\text{at degree }35,
\end{aligned}
\]

up to ordering. The standard-cycle theorem and power pullback make both
blocks algebraic. The length-eight block lives on a Fermat sixfold, only
its explicitly known standard cycle is needed, not the general Hodge
conjecture for that sixfold.

The exact identity is

\begin{equation}\label{eq:transfer210-5.1}
\boxed{
\begin{aligned}
 b\uplus P_{24}\uplus P_{27}\uplus P_{105}\uplus Q_3\uplus S_5\uplus S_7
 ={}&s\uplus P_{12}\uplus P_{30}\uplus P_{42}\uplus P_{102}\\
 &\uplus Q_9\uplus Q_{24}\uplus Q_{27}\uplus R_2.
\end{aligned}}\end{equation}

Both sides contain exactly 30 entries, with equal multiplicity of every
residue. Equivalently, \(\nu(b)-\nu(s)\) is the signed sum of the
fourteen displayed auxiliary generators, each with coefficient \(+1\) or
\(-1\). The certificate is stored in \texttt{results/transfer210.json}.

\hypertarget{transfer210-algebraicity-of-the-degree-210-packet}{%
\subsection{Algebraicity of the degree-210
packet}\label{transfer210-algebraicity-of-the-degree-210-packet}}

Assuming the inherited degree-70 carrier proposition, every rational
class in \(W_{[b]}\) is algebraic.

\textbf{Proof.} Let \(L\) be the auxiliary multiset to the left of \eqref{eq:transfer210-5.1}
and \(R\) the auxiliary multiset to its right. Thus

\[
b\uplus L=s\uplus R,
\qquad |L|=|R|=24.
\]

All component packets of \(L\) and \(R\) are algebraic. Negation
preserves their algebraicity, either by the rational Galois packet
statement or by conjugating the corresponding cycle class. Append \(-L\)
to both sides:

\begin{equation}\label{eq:transfer210-5.2}
b\uplus L\uplus(-L)=s\uplus R\uplus(-L).\end{equation}

The right-hand side is algebraic by concatenation, because the source
and every auxiliary component are algebraic. The left-hand auxiliary
part \(L\uplus(-L)\) is a multiset of exactly 24 complementary pairs.
The full character has 54 entries and is a character on
\(X_{210}^{52}\). Complementary-pair removal therefore gives
algebraicity of \(b\) on \(X_{210}^4\).

The character projector and its rational Galois orbit transfer this
statement to every class in \(W_{[b]}\). Its dimension is 48. This
proves the proposition. \hfill\ensuremath{\square}

For every \(k\ge1\), coordinate-power pullback gives the associated
packet on \(X_{210k}^4\). It does not establish the Hodge conjecture for
all packets on any of these fourfolds.

This proof does not subtract unknown algebraic cycles, assert
cancellation of arbitrary motives, or assume a cycle-existence theorem
for the new packet. Every negative coefficient in the incidence identity
is removed through explicitly algebraic auxiliary blocks and the
classical complementary-pair theorem.

\begin{figure}[tbp]\centering\includegraphics[width=0.98\linewidth]{24_transfer210_incidence.pdf}\caption{An exact signed multiplicity equality. Its algebraic implication also uses the stated classical correspondence operations.}\end{figure}

\hypertarget{transfer210-a-global-obstruction-to-the-direct-matrix-route}{%
\section{A global obstruction to the direct matrix
route}\label{transfer210-a-global-obstruction-to-the-direct-matrix-route}}

The previous search criterion required a nonsingular nonnegative
integral matrix \(A\), with columns \(v_j\), satisfying

\begin{equation}\label{eq:transfer210-6.1}
b\cdot v_j\equiv0\pmod {210},\qquad
w=A^{-1}\mathbf1>0,\qquad
\sum_jw_j>1.\end{equation}

The following exact certificate rules this out for the new target.

Set

\begin{equation}\label{eq:transfer210-6.2}
\eta=\frac1{420}(31,87,57,101,84,60).\end{equation}

Its coordinates are positive and sum to one. For every nonzero
\(v\in\mathbf Z_{\ge0}^6\) satisfying \(b\cdot v\equiv0\pmod {210}\),

\begin{equation}\label{eq:transfer210-6.3}
\boxed{\eta\cdot v\ge1.}\end{equation}

\hypertarget{transfer210-proof-of-the-universal-inequality}{%
\subsection{Proof of the universal
inequality}\label{transfer210-proof-of-the-universal-inequality}}

If \(|v|\ge14\), the numerator of \(\eta\cdot v\) is at least
\(14\cdot31=434>420\). It therefore suffices to check \(1\le |v|\le13\).

The exact verifier enumerates all 27,131 such nonnegative six-tuples.
Exactly 143 satisfy the invariant-column congruence. Each has numerator
at least 420, and the minimum is exactly 420. The enumerator uses
integer arithmetic only and includes all compositions in the specified
range. Together with the large-degree bound, this proves \eqref{eq:transfer210-6.3} for every
column, not just those within a search cutoff.

As an independent implementation, place residue states in
\(\mathbf Z/210\) and give the six transitions \(r\mapsto r+b_i\)
integer costs \((31,87,57,101,84,60)\). A nonempty zero-sum column is a
return walk. Exact shortest-path computation finds minimum nonempty
return cost 420. All costs are positive. This verifies the same
inequality by a different method.

\hypertarget{transfer210-no-direct-positive-weight-matrix-carrier}{%
\subsection{No direct positive-weight matrix
carrier}\label{transfer210-no-direct-positive-weight-matrix-carrier}}

If a matrix in \eqref{eq:transfer210-6.1} existed, then

\[
1=\eta^T\mathbf1=\eta^TAw
 =\sum_jw_j(\eta\cdot v_j)\ge\sum_jw_j.
\]

This contradicts \(\sum_jw_j>1\). \hfill\ensuremath{\square}

The single invariant column \(v=\mathbf1\), with coefficient one,
supplies a matching value-one witness for the relaxed column linear
program. Thus its global optimum is exactly one. This is not a claim
that a nonsingular six-column positive solution attains that optimum.

The invariant-column condition is unchanged by unit-conjugating \(b\),
permuting its coordinates with the corresponding dual permutation, or
inflating \(b\) and the modulus together. The obstruction persists under
those operations.

the algebraicity-transfer proposition above and the direct-matrix
obstruction corollary above together show a definite limitation of the
direct carrier criterion: a packet can be algebraic by a certified
transfer while every direct matrix of this type is excluded. This does
not rule out all geometric carriers or non-diagonal correspondences.

\begin{figure}[tbp]\centering\includegraphics[width=0.98\linewidth]{27_global_dual210.pdf}\caption{The exact global dual combines a complete finite premise with a uniform large-degree inequality.}\end{figure}

\hypertarget{transfer210-a-uniform-infinite-family-excluded-by-the-same-criterion}{%
\section{A uniform infinite family excluded by the same
criterion}\label{transfer210-a-uniform-infinite-family-excluded-by-the-same-criterion}}

For any integer \(q>5\), let

\begin{equation}\label{eq:transfer210-7.1}
a_q=(1,1+q,1+2q,1+3q,1+4q,5q-5)\pmod {5q}.\end{equation}

\hypertarget{transfer210-theorem-7.2}{%
\subsection{Theorem 7.2}\label{transfer210-theorem-7.2}}

Every \(a_q\) is a full-support middle-balanced character. No matrix
satisfying the previous direct positive-weight criterion exists for
\(a_q\).

\textbf{Proof of balance.} Let \(t\) be a unit modulo \(5q\), and let
\(s\) be its least positive residue modulo \(q\). The first five
multiplied residues permute

\[
s,s+q,s+2q,s+3q,s+4q.
\]

Their sum is \(5s+10q\). The last multiplied residue is \(5(q-s)\), so
the total is \(15q=3(5q)\).

\textbf{Proof of the obstruction.} Set

\[
\eta_q=\frac1q(1,1,1,1,1,q-5).
\]

This vector is positive and has sum one. For a nonzero invariant column
\(v\ge0\), put \(S=v_0+\cdots+v_4\). Reducing
\(a_q\cdot v\equiv0\pmod {5q}\) modulo \(q\) gives
\(S-5v_5\equiv0\pmod q\). Therefore

\[
\eta_q\cdot v=\frac{S+(q-5)v_5}{q}
=\frac{S-5v_5}{q}+v_5
\]

is a positive integer, and hence at least one. Apply the same matrix
inequality as in the preceding dual argument. \hfill\ensuremath{\square}

This is a proof for every \(q>5\), not an extrapolation from finite
runs. The verifier's checks for \(6\le q\le80\) are benchmarks only.
These packets are themselves the standard family \(\sigma_{5,1}\), so
the obstruction is to the method, not to algebraicity.

\hypertarget{transfer210-a-stronger-obstruction-holomorphic-powers-of-a-hodge-character}{%
\section{A stronger obstruction: holomorphic powers of a Hodge
character}\label{transfer210-a-stronger-obstruction-holomorphic-powers-of-a-hodge-character}}

\hypertarget{transfer210-proposition-8.1}{%
\subsection{Proposition 8.1}\label{transfer210-proposition-8.1}}

Let \(a\) be a full-support Fermat fourfold character. Suppose some
integer multiple

\[
c=ka\pmod m
\]

has all entries nonzero and least-positive residue sum \(m\). A direct
quotient by a subgroup of the diagonal automorphism group that captures
\(a\) cannot have a smooth projective uniruled compactification.

\textbf{Proof.} If \(K\) is a diagonal deck subgroup and
\(\chi_a|_K=1\), then

\[
\chi_c|_K=\chi_a^k|_K=1.
\]

The residue description places the \(c\)-line in \(H^{4,0}\) under the
direct-action convention (or its conjugate under the inverse-character
convention). Its rational Galois span contains a nonzero \((4,0)\) part
and is also \(K\)-invariant.

Full support makes restriction to the open coordinate torus injective on
this character space. As in the predecessor, each boundary stratum is
fixed by a coordinate scaling subgroup on which the character is
nontrivial, so the relevant Gysin terms vanish. Finite étale descent
gives the invariant pure weight-four space on the open quotient.

For any smooth projective compactification, the lowest-weight image
theorem and the semisimplicity of polarizable pure Hodge structures lift
this space from open cohomology to compact cohomology. In particular the
compactification has a nonzero \(H^{4,0}\). These are the same classical
mixed-Hodge inputs as in the predecessor, see \cite{ArapuraUniruled},
\cite{DeligneHodgeII}.

A smooth projective uniruled fourfold has no nonzero holomorphic top
form: a generically finite dominant parametrization from a resolution of
a threefold times \(\mathbf P^1\) would pull such a form back
injectively, whereas the latter has none. This is a contradiction. \hfill\ensuremath{\square}

The assertion concerns this direct diagonal-quotient construction, not
arbitrary algebraic correspondences. It persists after coordinate-power
inflation of the source.

\Needspace{10\baselineskip}

\hypertarget{transfer210-exact-witnesses}{%
\subsection{Exact witnesses}\label{transfer210-exact-witnesses}}

{\small\begin{longtable}[]{@{}
  >{\raggedright\arraybackslash}p{(\columnwidth - 6\tabcolsep) * \real{0.2308}}
  >{\raggedright\arraybackslash}p{(\columnwidth - 6\tabcolsep) * \real{0.2308}}
  >{\raggedleft\arraybackslash}p{(\columnwidth - 6\tabcolsep) * \real{0.3077}}
  >{\raggedright\arraybackslash}p{(\columnwidth - 6\tabcolsep) * \real{0.2308}}@{}}
\toprule\noalign{}
\begin{minipage}[b]{\linewidth}\raggedright
Modulus
\end{minipage} & \begin{minipage}[b]{\linewidth}\raggedright
Middle-balanced character a
\end{minipage} & \begin{minipage}[b]{\linewidth}\raggedleft
k
\end{minipage} & \begin{minipage}[b]{\linewidth}\raggedright
Full-support character ka with sum m
\end{minipage} \\
\midrule\noalign{}
\endhead
\bottomrule\noalign{}
\endlastfoot
110 & (1,24,62,71,81,91) & 11 & (11,44,22,11,11,11) \\
110 & (1,31,55,71,81,91) & 11 & (11,11,55,11,11,11) \\
114 & (1,13,43,72,103,110) & 27 & (27,9,21,6,45,6) \\
114 & (1,13,43,80,102,103) & 9 & (9,3,45,36,6,15) \\
114 & (1,7,78,79,86,91) & 3 & (3,21,6,9,30,45) \\
220 & (1,62,111,142,162,182) & 11 & (11,22,121,22,22,22) \\
\end{longtable}}

Each source's middle balance and each product's nonzero entries and
residue sum are checked exactly. The multipliers are nonunits, so there
is no contradiction with the middle balance of all \textbf{unit}
conjugates of the source character.

These witnesses explain why merely increasing the direct diagonal
carrier search cannot address these targets. They do not show that the
targets are nonalgebraic.

\hypertarget{transfer210-computation-contract-failures-and-remaining-work}{%
\section{Computation contract, failures and remaining
work}\label{transfer210-computation-contract-failures-and-remaining-work}}

\hypertarget{transfer210-exact-final-replay}{%
\subsection{Exact final replay}\label{transfer210-exact-final-replay}}

\texttt{python\ checks/verify.py} needs Python 3.10 or later and only
its standard library. It runs without network access, optimization
software or floating-point arithmetic. It verifies the frozen identity,
all named character conditions, the standard-block deflations, the
54-entry stabilization, the dual's complete finite premise and degree
cutoff, an independent exact shortest-cycle check, six holomorphic-power
witnesses and 75 infinite-family benchmarks.

The code does not construct equations for a basis of representing
surfaces, verify the external theorems, or compile a proof in Lean. The
standard-block and carrier constructions give existence through known
algebraic correspondences, not a list of all representing surfaces.

\hypertarget{transfer210-discovery-distinct-from-proof}{%
\subsection{Discovery, distinct from
proof}\label{transfer210-discovery-distinct-from-proof}}

The transfer was found by a bounded mixed-integer search among
complementary pairs and standard blocks at degree 210. The solver uses
floating-point internals. The stored result has fourteen nonzero
coefficients and was accepted only after exact integer replay. Its
reported optimizer optimum is not a claim that fourteen is globally
minimal among all possible geometric identities.

An earlier Hermite-normal-form membership attempt exceeded its 20-second
execution allowance and was stopped without a conclusion. Direct
integer-programming discovery replaced it successfully. Bounded
degree-18 matrix scans were exploratory and never used as global
impossibility proofs. The exact dual arguments in the two dual
obstruction arguments above replace those bounded observations for the
targets they cover.

Original publisher PDF and some ancillary-log fetches failed. No missing
log was used as evidence, the transfer and its verification were
produced locally. The source-access limit for the classical Aoki theorem
is resolved in the source-access discussion above.

\hypertarget{transfer210-what-remains}{%
\subsection{What remains}\label{transfer210-what-remains}}

The general Hodge conjecture is not established. This block does not
cover all packets on \(X_{210}^4\), establish universal closure of the
Fermat packet calculus, or reduce arbitrary projective varieties to
diagonal hypersurfaces. It does not complete the earlier central-defect
realization lemma or a general CM half-twist Chow correspondence.

The direct matrix criterion cannot be the universal bridge: the two dual
obstruction arguments above disprove that assertion. Direct diagonal
uniruled quotients cannot capture the holomorphic-power examples either.
A further positive argument must therefore use stabilization and
transfer, a different known-algebraic carrier class, or genuinely
non-diagonal correspondences. That is a mathematical constraint on the
next construction, not a conjectured failure of Hodge.
